% ============================================================================
%  Precautionary Saving in Europe --- Feedback Implementation
%  Liquidity, term structure, risk premia, and a forward-looking lens.
%  Compile with:  pdflatex documentation.tex   (run twice for references)
%  Font: Palatino (mathpazo). Matches the main project documentation style.
% ============================================================================
\documentclass[11pt,a4paper]{article}

% ---- Fonts: Palatino for text and math -------------------------------------
\usepackage[T1]{fontenc}
\usepackage[utf8]{inputenc}
\usepackage{mathpazo}
\linespread{1.05}
\usepackage{microtype}

% ---- Layout ----------------------------------------------------------------
\usepackage[a4paper,margin=2.6cm]{geometry}
\usepackage{booktabs}
\usepackage{array}
\usepackage{amsmath,amssymb}
\usepackage{graphicx}
\graphicspath{{figures/}{./}}
\usepackage{caption}
\captionsetup{font=small,labelfont=bf}
\usepackage{enumitem}
\setlist{itemsep=2pt,topsep=3pt}
\usepackage[section]{placeins}   % keep each section's floats within it (no drift)
% pack pages tighter (fewer float-only pages / less trailing whitespace)
\renewcommand{\topfraction}{0.92}
\renewcommand{\bottomfraction}{0.85}
\renewcommand{\textfraction}{0.08}
\renewcommand{\floatpagefraction}{0.75}
\usepackage{xcolor}
\definecolor{codebg}{gray}{0.96}
\definecolor{rule}{gray}{0.5}

\usepackage[hidelinks,colorlinks=true,linkcolor=black!70!blue,
            urlcolor=black!60!blue,citecolor=black]{hyperref}

% ---- Convenience macros ----------------------------------------------------
\newcommand{\code}[1]{\texttt{#1}}
\newcommand{\dataset}[1]{\textsc{#1}}
% Robust figure inclusion: shows a placeholder if the PNG is absent.
\newcommand{\fig}[3]{% {width-fraction}{file-without-ext}{caption}
  \begin{figure}[htbp]\centering
    \IfFileExists{figures/#2.png}
      {\includegraphics[width=#1\linewidth]{#2}}
      {\fbox{\parbox{0.80\linewidth}{\centering\itshape figure not found ---
        run \code{bash run.sh} (this one needs FRED) to generate it.}}}
    \caption{#3}\label{fig:#2}
  \end{figure}}

\title{\textbf{Precautionary Saving in Europe: Feedback Implementation}\\[3pt]
       \large Liquidity, Term Structure, Risk Premia, and a Forward-Looking Lens}
\author{Project documentation --- feedback extension}
\date{\today}

\begin{document}
\maketitle
\thispagestyle{empty}

\begin{abstract}
\noindent
This addendum implements supervisor feedback on the savings analysis. Its
organising idea is a single methodological correction: stop treating
\emph{cash and deposits} as the only precautionary buffer, and instead rank
household assets on a \emph{liquidity / maturity ladder}. On that ladder, the
post-2022 move \emph{out of} instant-access cash --- which the project's
``follow-the-money'' extension had read as evidence of \emph{yield-chasing
rather than precaution} --- looks very different: the money did not leave the
liquid buffer, it moved up the risk/return scale while staying sellable within
days. We then add the four further pieces the feedback asked for: the money
supply divided by term ($M1$/$M2$/$M3$); observable risk premia set against
geopolitical tension; a forward-looking expectations composite; and an
energy-price liquidity-adequacy analysis; debt securities are split by term for a
finer ladder, and a United States comparison (Fed Z.1) sets the euro-area picture
in relief. All series are free (Eurostat, the ECB Data Portal, FRED, the
Caldara--Iacoviello GPR index, and the US Financial Accounts); the scripts reuse
the main project's helpers. Headline result: the broad
sellable-fast share of household saving flows \emph{rose} from $54\%$ to $73\%$
after 2022, and about half of household financial wealth is sellable within
days --- so the precaution-versus-yield dichotomy is false. Households'
forward-looking income expectations predict the saving rate far better than the
geopolitical index does, and the buffer is concentrated up the income
distribution, leaving the most energy-exposed households uncushioned. Figures
come from a clean live run on \today; they are revised by the providers and
should be regenerated before citation.
\end{abstract}

\tableofcontents
\bigskip
\hrule
\bigskip

% ============================================================================
\section{The question, revisited}\label{sec:question}
% ============================================================================

\paragraph{Where the main project left off.} The project tests whether the euro
area's elevated household saving rate is \emph{substantially precautionary} ---
a response to uncertainty rather than to income, wealth, or the interest rate
alone. The time-series leg was supportive; the cross-country and formal
identification legs were inconclusive. The sharpest negative result came from
the ``follow-the-money'' extension: after 2022 the share of saving going into
instant-access cash collapsed while money flowed into bonds and term deposits.
Read through a binary --- \emph{instant-access cash $=$ precautionary},
\emph{bonds and term deposits $=$ yield-chasing} --- this looked like evidence
\emph{against} precaution.

\paragraph{The feedback.} The supervisors' objection targets exactly that
binary. Many assets can be sold fast, so cash and deposits are not the only
precautionary store of value; we should look harder at each asset type and at
the money supply, \emph{divided by term}; examine risk premia and whether they
rise with geopolitical tension; adopt a more forward-looking approach; and
discuss the risk of holding too little cash to absorb rising energy prices ---
in short, \emph{when} households save, \emph{where}, and how \emph{liquid} those
holdings are.

\paragraph{The organising reframe.} A bond, a listed share, or a money-market
fund unit is sellable in days. Treating it as ``locked'' understates the
precautionary buffer. We therefore replace the binary with a four-rung
liquidity / maturity ladder (Table~\ref{tab:ladder}); we divide finer where the
data allow --- debt securities are split \emph{by term}, short-term (F31) into
near-money and long-term (F32) into the marketable tier --- and recompute the
composition of saving on it. The monetary aggregates are the central bank's own
liquidity-by-term tiering of the money-holding sector's claims, so they are the
natural macro counterpart. The remaining sections add the risk-premium,
forward-looking, and energy dimensions, and Section~\ref{sec:synthesis} states
what all of this implies for the precautionary question.

% ============================================================================
\section{Data and method}\label{sec:data}
% ============================================================================

Every script is standalone, reuses the shared helpers in \code{\_common.py}
(lifted from the main project's modules), and pulls only free, keyless data,
writing figures to \code{figures/} and a report plus CSVs to \code{data/}.
Table~\ref{tab:data} lists the additional series. Two helpers are new: a
generalised \code{household\_instruments()} that runs the \emph{same} tier
classification on flows (\dataset{nasa\_10\_f\_tr}) and on stocks
(\dataset{nasa\_10\_f\_bs}); and \code{ecb\_sdmx()}, a keyless pull from the ECB
Data Portal for the monetary aggregates.

\begin{table}[htbp]\centering\small
\caption{The liquidity / maturity ladder (ESA 2010 instrument codes).}
\label{tab:ladder}
\begin{tabular}{@{}p{3.4cm}p{5.4cm}p{6.0cm}@{}}
\toprule
\textbf{Tier} & \textbf{Instruments} & \textbf{Rationale} \\
\midrule
T1 --- instant ($\approx M1$) & F21 currency, F22 overnight deposits &
  settle immediately, no price risk \\
T2 --- near-money ($\approx M2/M3$) & F29 other deposits (term \& redeemable at
  notice), F521 MMF shares, \textbf{F31 short-term debt securities} &
  accessible at short notice / maturity, little duration risk \\
T3 --- marketable & \textbf{F32 long-term debt securities}, F511 listed shares,
  F522 non-MMF investment-fund shares &
  sellable within days at market price (duration / price risk) \\
T4 --- illiquid / contractual & F512, F519 unlisted \& other equity; F6
  insurance, pension \& guarantee schemes &
  locked, contractual, or with no ready secondary market \\
\bottomrule
\end{tabular}
\end{table}

\begin{table}[htbp]\centering\small
\caption{Additional data series (all free / keyless).}
\label{tab:data}
\begin{tabular}{@{}p{4.6cm}p{5.0cm}p{5.2cm}@{}}
\toprule
\textbf{Series} & \textbf{Source} & \textbf{Used in} \\
\midrule
Household financial balance sheet (stocks) & Eurostat \dataset{nasa\_10\_f\_bs} &
  \S\ref{sec:ladder} \\
Household financial transactions (flows) & Eurostat \dataset{nasa\_10\_f\_tr} &
  \S\ref{sec:ladder} \\
Monetary aggregates $M1,M2,M3$ & ECB Data Portal (\dataset{BSI}); FRED fallback &
  \S\ref{sec:money} \\
Euro high-yield credit spread (OAS) & FRED \code{BAMLHE00EHYIOAS} &
  \S\ref{sec:premia} \\
Italy \& Germany 10y yields & FRED \code{IRLTLT01ITM156N}, \code{IRLTLT01DEM156N} &
  \S\ref{sec:premia} \\
Implied equity volatility (VIX) & FRED \code{VIXCLS} &
  \S\ref{sec:premia}, \S\ref{sec:fwd} \\
Consumer survey (saving, unemployment) & Eurostat \dataset{ei\_bsco\_m} &
  \S\ref{sec:fwd} \\
Energy HICP & Eurostat \dataset{prc\_hicp\_midx} / \dataset{manr} (NRG) &
  \S\ref{sec:energy} \\
\bottomrule
\end{tabular}
\end{table}

A guard worth stating up front: Eurostat does not always publish the F5
sub-instruments (F511/F512/F519/F521/F522). The classifier auto-detects what is
present and falls back \emph{conservatively} --- an un-splittable equity/funds
block is assigned to the illiquid tier T4 --- so every ``broad-liquid'' figure
below is a \emph{lower bound}. We never overstate liquidity. In the run reported
here the MMF split (F521) was unavailable, so fund shares enter T3 whole.

% ============================================================================
\section{The liquidity ladder}\label{sec:ladder}
% ============================================================================

\paragraph{Stocks: what households hold.} Re-ranking the household financial
balance sheet (Figure~\ref{fig:liquidity_ladder_stocks},
Table~\ref{tab:stocks}), instant cash (T1) is only about a fifth of financial
wealth, but the \emph{broad sellable-fast} share (T1$+$T2$+$T3) is around
\emph{half}. The remainder (T4 $\approx 47\%$) is dominated by insurance and
pension claims and by unlisted business equity --- genuinely illiquid. The point
the feedback asked for is immediate: the mobilisable buffer is far larger than
cash.

\fig{0.92}{liquidity_ladder_stocks}{Euro-area household financial wealth by
liquidity tier (stocks, \dataset{nasa\_10\_f\_bs}). Only $\sim18\%$ is instant
cash (T1); the dashed line --- the broad sellable-fast share T1$+$T2$+$T3 --- sits
near $50\%$.}

\begin{table}[htbp]\centering\small
\caption{Stock shares of household financial wealth (\%), pre- vs post-shock
averages.}\label{tab:stocks}
\begin{tabular}{@{}lrrr@{}}
\toprule
 & 2015--19 & 2022+ & change (pp) \\
\midrule
T1 instant (cash, overnight)        & 17.2 & 19.2 & $+1.9$ \\
T2 near-money (term/notice, MMF)    & 14.5 & 13.0 & $-1.5$ \\
T3 marketable (bonds, shares, funds)& 15.8 & 17.6 & $+1.9$ \\
T4 illiquid (unlisted eq., ins./pension) & 49.1 & 46.8 & $-2.3$ \\
\midrule
\emph{narrow cash (T1)}             & 17.2 & 19.2 & $+1.9$ \\
\emph{broad sellable-fast (T1$+$T2$+$T3)} & 47.5 & 49.8 & $+2.3$ \\
\bottomrule
\end{tabular}
\end{table}

\paragraph{Flows: where the marginal saving went.} The reframe is starkest in
the flows (Figure~\ref{fig:liquidity_ladder_flows}, Table~\ref{tab:flows}). The
share of yearly financial saving going into instant cash did collapse, from
$66\%$ (2015--19) to $2\%$ (2022+) --- the headline of the old binary. But the
\emph{broad sellable-fast} share \emph{rose}, from $54\%$ to $73\%$: the money
left cash and went into term deposits, money-market funds, bonds and listed
shares, every one of which can be sold within days. On the ladder this is not
``leaving the precautionary buffer''; it is holding the buffer in a
higher-yielding but still-liquid form once positive rates made cash costly to
hold.

\fig{0.92}{liquidity_ladder_flows}{Euro-area household net financial flows by
liquidity tier (\dataset{nasa\_10\_f\_tr}). The instant-cash line collapses after
2022 while the marketable tier (T3) surges --- a move \emph{within} the liquid
part of the ladder.}

\begin{table}[htbp]\centering\small
\caption{Flow shares of yearly financial saving (\%), pre- vs post-shock
averages. (Net-flow shares can exceed 100\% or go negative; read the levels in
Figure~\ref{fig:liquidity_ladder_flows} alongside.)}\label{tab:flows}
\begin{tabular}{@{}lrrr@{}}
\toprule
 & 2015--19 & 2022+ & change (pp) \\
\midrule
T1 instant (cash, overnight)        & 66.2 & 2.3 & $-63.9$ \\
T2 near-money (term/notice, MMF)    & $-10.7$ & 33.4 & $+44.2$ \\
T3 marketable (bonds, shares, funds)& $-1.0$ & 37.4 & $+38.4$ \\
T4 illiquid (unlisted eq., ins./pension) & 43.0 & 20.7 & $-22.3$ \\
\midrule
\emph{narrow cash (T1)}             & 66.2 & 2.3 & $-63.9$ \\
\emph{broad sellable-fast (T1$+$T2$+$T3)} & 54.4 & 73.0 & $+18.6$ \\
\bottomrule
\end{tabular}
\end{table}

This is the single most important result for the precautionary question: it
removes the strongest piece of evidence \emph{against} the hypothesis. ``Out of
cash'' is not ``out of the buffer.'' It does not, by itself, \emph{prove} the
motive is precaution --- staying liquid is consistent with precaution, not proof
of it --- but it restores precaution as a live reading after the binary had
seemingly ruled it out.

% ============================================================================
\section{Money supply, divided by term}\label{sec:money}
% ============================================================================

The monetary aggregates tell the same story at the macro level
(Figure~\ref{fig:money_supply_by_term}, Table~\ref{tab:money}). As the ECB
raised rates, money migrated \emph{out of} instant overnight balances ($M1$ fell
by about EUR\,1.2tn between 2022 and 2024) and \emph{into} term and notice
deposits ($M2-M1$ rose from EUR\,3.4tn to EUR\,4.9tn) and marketable near-money
($M3-M2$: MMF shares, repo, short debt). Broad money ($M3$) kept growing and
co-moves with the household saving rate (annual correlation $+0.27$). Money
itself is tiered by term, and the buffer is not just $M1$ cash.

\fig{0.86}{money_supply_by_term}{Euro-area money supply decomposed by term:
$M1$ (instant), $M2-M1$ (term/notice deposits), $M3-M2$ (marketable). The
post-2022 widening of the $M2-M1$ band as $M1$ dips is the term shift.}

\begin{table}[htbp]\centering\small
\caption{Euro-area money supply by term (EUR\,bn, annual mean; ECB
\dataset{BSI}).}\label{tab:money}
\begin{tabular}{@{}lrrrr@{}}
\toprule
year & $M1$ instant & $M2-M1$ term/notice & $M3-M2$ marketable & $M3$ \\
\midrule
2019 &  8{,}659 & 3{,}447 &   640 & 12{,}745 \\
2021 & 10{,}846 & 3{,}434 &   766 & 15{,}045 \\
2022 & 11{,}556 & 3{,}604 &   776 & 15{,}936 \\
2023 & 10{,}720 & 4{,}391 &   900 & 16{,}011 \\
2024 & 10{,}351 & 4{,}938 & 1{,}096 & 16{,}384 \\
\bottomrule
\end{tabular}
\end{table}

\paragraph{Honest reading.} The $M1\!\rightarrow\!$ term migration is exactly
what rising rates predict on their own, so this corroborates ``the buffer
relocated, it did not disappear'' but cannot by itself separate precaution from
intertemporal substitution.

% ============================================================================
\section{Risk premia and geopolitical tension}\label{sec:premia}
% ============================================================================

The feedback asked whether risk premia rise with geopolitical tension, and
whether we can observe it. Three euro-area premia are free, observable, and can
be set against the Geopolitical Risk (GPR) index the project already pulls: the
\emph{credit} premium (euro high-yield option-adjusted spread); the
\emph{sovereign / fragmentation} premium (the Italy--Germany 10-year spread, the
cleanest euro-area gauge); and \emph{risk aversion} (the VIX, with VSTOXX as the
euro analogue). The script overlays each on GPR, reports the co-movement in
levels and monthly changes, and event-windows the February 2022 invasion.

\fig{0.92}{risk_premia_vs_gpr}{Euro-area credit and sovereign spreads (top) and
VIX (bottom) against the GPR index. Premia rise and jump with geopolitical
tension.}

\paragraph{Status.} The figure shows the \emph{sovereign} (Italy--Germany) spread,
pulled from the ECB Data Portal --- a fallback that renders this gauge even when
FRED is unreachable. The BTP--Bund spread jumped about $+0.6$\,pp at the February
2022 invasion; over the full sample its \emph{level} correlation with GPR is muted
because the 2011--12 euro debt crisis dominates the series, so the event-window
jump is the cleaner read. Running \code{python3 risk\_premia.py} with FRED access
adds the credit spread (euro high-yield OAS) and the VIX, which show clear positive
GPR co-movement and Feb-2022 jumps. Either way the point holds: geopolitical
tension is observably priced --- an \emph{enabling} fact (the uncertainty is real
and measurable in markets) rather than direct evidence on the household saving
motive.

% ============================================================================
\section{A forward-looking lens}\label{sec:fwd}
% ============================================================================

The core analysis relates the \emph{realised} saving rate to
\emph{contemporaneous} uncertainty. A forward-looking version asks whether what
households and markets \emph{expect} carries information about saving. We build a
caution composite from variables that are, by construction, about the future ---
expected unemployment and intended saving (Eurostat \dataset{ei\_bsco\_m}), and
implied volatility (VIX) when FRED is reachable --- as the mean of their
$z$-scores, and test it against the realised saving rate
(Figures~\ref{fig:forward_looking_composite}--\ref{fig:forward_looking_leadlag}).

\fig{0.86}{forward_looking_composite}{The forward-looking caution composite
($z$) against the realised saving rate ($z$).}

\fig{0.62}{forward_looking_leadlag}{Lead--lag correlation
$\mathrm{corr}(\text{saving}_t,\ \text{composite}_{t-k})$; $k>0$ means the
composite leads.}

\paragraph{Result.} A one-quarter-ahead regression,
$\text{saving}_t = \alpha + \beta\,\text{composite}_{t-1}$, gives $\beta=+1.88$
($p<0.001$, $R^2=0.38$, $n=107$), whereas the \emph{contemporaneous} GPR
benchmark explains essentially nothing ($R^2\approx0.00$). Two honest caveats.
First, the lead--lag correlation peaks \emph{contemporaneously} ($+0.76$ at
$k=0$), so the right claim is ``an already-observable composite predicts saving
far better than GPR,'' not ``expectations strongly lead saving.'' Second, the
composite includes \emph{intended} saving, which is close to the target by
construction; the cleaner forward signal is \emph{expected unemployment} (and,
with FRED, the VIX). The substantive implication for the question is a
refinement: the operative uncertainty looks like households' own perceived
\emph{income / labour-market} risk, \emph{not} the abstract geopolitical index
--- GPR barely moves household saving.

% ============================================================================
\section{Energy prices and the adequacy of cash}\label{sec:energy}
% ============================================================================

A precautionary buffer is only useful if it can cover a \emph{non-deferrable}
shock; the energy bill is the cleanest example. Two findings
(Figures~\ref{fig:energy_vs_cash_buffer}--\ref{fig:energy_squeeze_by_income}).
First, the euro-area energy price level rose about $60\%$ above its 2019 level
and stayed elevated, while the instant-cash share of household wealth did
\emph{not} rise to meet it (it touched $21.5\%$ in 2022, then fell back to
$\sim18\%$ by 2024) --- the non-deferrable bill grew as the cash cushion thinned,
and households were moving \emph{out} of cash (\S\ref{sec:ladder}).

\fig{0.86}{energy_vs_cash_buffer}{The euro-area energy price level (left, 2019$=$100)
against the instant-cash share of household wealth (right).}

Second, the squeeze is \emph{regressive} (Table~\ref{tab:energy}): the
lowest-income quintile spends the largest budget share on energy yet
\emph{dissaves} (a negative saving rate), while the top quintile saves heavily on
the smallest energy share. The aggregate ``precautionary buffer'' is therefore a
top-of-distribution phenomenon, and the households most exposed to an energy
shock are the least able to self-insure --- they have no buffer to draw on, cash
or otherwise.

\fig{0.78}{energy_squeeze_by_income}{Energy budget share (illustrative) versus
the saving rate, by income quintile. Highest burden meets the thinnest buffer.}

\begin{table}[htbp]\centering\small
\caption{Energy burden and saving capacity by income quintile (Q1 $=$ lowest).
Energy shares are illustrative/stylised --- verify against the Household Budget
Survey; saving rates are from Eurostat experimental ICW.}\label{tab:energy}
\begin{tabular}{@{}lrrrrr@{}}
\toprule
 & Q1 & Q2 & Q3 & Q4 & Q5 \\
\midrule
energy share of spending (\%) & 11.0 & 9.5 & 8.0 & 7.0 & 5.5 \\
saving rate (\%) & $-3.4$ & $+17.6$ & $+26.0$ & $+33.9$ & $+44.2$ \\
\bottomrule
\end{tabular}
\end{table}

\paragraph{A caveat to the reframe.} ``Sellable fast'' is not ``safe to spend
now.'' Liquidating marketable (T3) assets to pay an energy bill in a stressed,
high-tension state crystallises losses precisely when those asset prices are
depressed and risk premia (\S\ref{sec:premia}) are elevated. Liquidity in normal
times is not the same as adequacy in a crisis --- and the most vulnerable
households cannot do even this.

% ============================================================================
\section{United States comparison}\label{sec:us}
% ============================================================================

Does the same picture hold across the Atlantic? We rebuild the liquidity ladder
for US households from the Federal Reserve Financial Accounts (Z.1) and set it
against the euro area (Figure~\ref{fig:us_vs_ea_ladder}), and contrast the saving
rates (Figure~\ref{fig:us_vs_ea_saving}). The expected contrast: US households
hold far \emph{less} in cash and deposits (T1/T2) and far \emph{more} in
marketable equities and investment-fund shares (T3) than euro-area households.
Their precautionary buffer is therefore even more ``sellable-fast'' --- but also
more exposed to \emph{market price risk}, which sharpens the
Section~\ref{sec:energy} caveat: a buffer held in equities can fall in value
exactly when it is needed. The US personal saving rate is also structurally lower
and more volatile than the euro-area household rate.

\fig{0.84}{us_vs_ea_ladder}{Household financial wealth by liquidity tier, US vs
euro area (stocks, latest). US: Fed Z.1 (households \& nonprofits); EA: Eurostat.}

\fig{0.84}{us_vs_ea_saving}{Saving rates: US (personal, \emph{net}) vs euro area
(household, \emph{gross}) --- compare the dynamics, not the levels.}

\paragraph{Status.} \code{us\_comparison.py} is FRED-sourced (Z.1 + the US saving
rate), so the two figures regenerate when it is run where FRED is reachable. The
US Z.1 series ids are best-effort and the script prints which resolved; shares are
taken against \emph{total} financial assets, so an unresolved component shows as a
residual rather than inflating the liquid share.

% ============================================================================
\section{What this says about the question}\label{sec:synthesis}
% ============================================================================

\begin{enumerate}[label=\arabic*.]
  \item \textbf{The false binary dissolves.} The post-2022 saving is best read as
  a \emph{liquid precautionary buffer that also earns yield}. The
  ``precaution versus yield-chasing'' dichotomy was an artefact of treating only
  cash as liquid; on the ladder the broad sellable-fast share of saving flows
  \emph{rose} (54\%\,$\to$\,73\%), and about half of household wealth is sellable
  within days (\S\ref{sec:ladder}, \S\ref{sec:money}).
  \item \textbf{The operative uncertainty is income fear, not geopolitics.}
  Households' forward-looking expectations predict the saving rate
  ($R^2\approx0.38$ one quarter ahead) while the geopolitical index does not
  ($R^2\approx0$) --- so the precautionary channel runs through perceived
  job/income risk (\S\ref{sec:fwd}). Risk premia confirm that tension is priced
  (\S\ref{sec:premia}), but that is an enabling fact, not the household motive.
  \item \textbf{The buffer is unevenly held.} The aggregate masks that the most
  energy-exposed, most income-fearful households have no cushion at all
  (\S\ref{sec:energy}).
\end{enumerate}

The verdict therefore moves from the presented version's cautious ``probably
yield-chasing'' to a sharper and better-defended statement: \emph{a liquid
precautionary buffer, optimised for yield once rates turned positive, driven more
by income fear than by geopolitics, and concentrated among the households who can
afford it.}

% ============================================================================
\section{Limitations}\label{sec:limits}
% ============================================================================
\begin{itemize}
  \item \textbf{Correlation, not causation.} Rates, inflation, and uncertainty
  all moved together after 2022; none of these descriptive results isolates the
  precautionary channel, and the main project's formal cross-country/IRF tests
  still return nulls.
  \item \textbf{Rate-consistency.} The money-supply term shift (\S\ref{sec:money})
  is equally consistent with a pure interest-rate story.
  \item \textbf{Mechanical component.} Part of the forward-looking $R^2$ reflects
  that the composite includes intended saving (\S\ref{sec:fwd}).
  \item \textbf{Lower-bound liquidity.} Where Eurostat does not split F5, the
  equity/funds block is treated as illiquid, so the broad-liquid shares are lower
  bounds (\S\ref{sec:data}).
  \item \textbf{Illustrative energy shares.} The by-income energy budget shares
  are stylised and should be verified against the Household Budget Survey.
  \item \textbf{Pending data.} The two US-comparison figures (and the credit /
  VIX series within the risk-premia figure) regenerate on a machine with FRED
  access; the risk-premia sovereign spread already renders via the ECB fallback.
  All series are pulled live and revised --- regenerate before citing.
\end{itemize}

\end{document}
